\documentclass[11pt]{article}

\usepackage[margin=1in]{geometry}
\usepackage{amsmath,amssymb}
\usepackage{siunitx}
\usepackage{booktabs}
\usepackage{xcolor}
\usepackage[colorlinks=true,linkcolor=blue,urlcolor=blue]{hyperref}

\newcommand{\En}{E_{\perp}}
\newcommand{\nhat}{\hat{\mathbf{n}}}
\newcommand{\code}[1]{\texttt{#1}}

\title{The no-flux (Neumann) boundary condition for the FR4 laminate\\
       in the \textsc{pochoir} pixel drift solver:\\
       why it is physically correct and why the FR4 surface-charge\\
       density need not appear as an explicit source term}
\author{\textsc{pochoir} --- insulating-surface-boundary branch}
\date{\today}

\begin{document}
\maketitle

\begin{abstract}
The drift-path overshoot at the inter-pixel gap of a LArTPC pixel plane is removed
by treating the FR4 laminate below the copper pads as a \emph{no-flux (homogeneous
Neumann)} boundary, $\partial\phi/\partial n = 0$, rather than as a dielectric
permittivity slab. This note argues that the Neumann condition is the correct
\emph{steady-state} model of a charged-up bulk insulator in a DC drift field, and
shows that the accumulated surface charge on the FR4 is not neglected but rather
\emph{absorbed into} the boundary condition: it is the response that enforces
$\En=0$, so the field can be solved without it as an input. The dielectric
(displacement) formulation is shown to be the wrong regime, consistent with the
observed worsening of the overshoot when a finite permittivity is used.
\end{abstract}

\section{Geometry and the discrete operator}

The pixel plane is a physical laminate: a \SI{1}{oz} copper pad
($\approx\SI{0.0348}{mm}$, one cell at \SI{0.05}{mm}) sitting on top of a
\SI{1.6}{mm} FR4 slab (32 cells). Electrons drift downward through liquid argon
(LAr) under a uniform field ($-\SI{50}{V/mm}$) and are collected on the pads at
$z=\SI{10}{mm}$; the FR4 occupies $z\in[8.4,10.0]\,\si{mm}$ directly beneath and
between the pads.

The solver relaxes the electrostatic potential $\phi$ with a masked
mirror/ghost-cell stencil. For a live cell adjacent to $k$ FR4 (``solid'')
neighbours, $\phi$ is updated as the average over only its $(2N-k)$ non-solid
neighbours, with the normalisation reduced to match:
\begin{equation}
  \phi_{\mathrm{new}}
     \;=\;
     \frac{1}{\,2N-k\,}\!\!\sum_{\substack{\text{neighbours }m \\ m\notin\text{FR4}}}\!\!\phi_m ,
  \label{eq:stencil}
\end{equation}
where $N$ is the spatial dimension. Dropping the FR4 bonds from \emph{both} the sum
and the denominator is the standard ghost-cell discretisation of a zero normal
derivative on every face of the excluded region,
\begin{equation}
  \left.\frac{\partial \phi}{\partial n}\right|_{\text{FR4 face}} = 0
  \qquad\Longleftrightarrow\qquad
  \En \equiv \mathbf{E}\cdot\nhat = 0 .
  \label{eq:neumann}
\end{equation}
No permittivity and no charge/source term enter Eq.~\eqref{eq:stencil}; the FR4
cells are frozen and excluded from the solve. In the code this is
\code{neumann\_coeff}/\code{stencil\_poisson\_neumann}
(\code{pochoir/fdm\_generic.py}), routed to whenever an insulator mask is supplied
(\code{pochoir/fdm\_torch.py}), with any supplied \code{epsilon} explicitly nulled.
On the drift side the velocity is set to zero inside the FR4 and a terminal event
stops a path at the FR4 top face (tagged \code{DRIFT\_SURFACE}); that tag is a
bookkeeping label, not a term in the field equation.

\section{Why $\En = 0$ is the correct physics}

\subsection{FR4 carries no free current}
FR4 is a bulk insulator with volume resistivity
$\rho \sim \SI{e14}{\ohm\,m}$. In DC steady state the conduction current density is
$\mathbf{J}=\sigma\mathbf{E}$ with $\sigma\to 0$, so the normal current into the
laminate surface vanishes,
\begin{equation}
  J_\perp = \sigma\,\En \approx 0 .
\end{equation}
Physically, a drifting electron in the LAr cannot enter or cross the FR4: it faces
a wide-bandgap insulator with no available conduction states. The interface is
impermeable to the drift current.

\subsection{The surface charges up until the normal field vanishes}
Because charge cannot pass \emph{through} the FR4, any field line that would
otherwise terminate on the laminate deposits electrons \emph{on its surface}. Those
accumulated electrons repel subsequently arriving electrons. This continues until,
in equilibrium, no further charge can land --- i.e.\ until the normal field at the
LAr/FR4 interface is driven to zero:
\begin{equation}
  \En\big|_{\text{LAr side}} \;\xrightarrow{\;\text{steady state}\;}\; 0 .
\end{equation}
A detector operates in exactly this charged-up steady state. The equilibrium
condition is therefore the homogeneous Neumann condition of
Eq.~\eqref{eq:neumann}. Its consequence is that the field lines run \emph{tangential}
to the FR4 and can only terminate on the conductors: a gap electron is
\emph{funnelled along the surface onto the nearest pad} instead of diving through
the laminate. This is precisely the mechanism that removes the gap overshoot.

\section{Why the FR4 surface-charge density is not an explicit term}

This is the subtle point: the surface charge is \emph{not neglected}; it is
\emph{absorbed into} the boundary condition.

\subsection{Free (accumulated) surface charge is the response, not an input}
The accumulated surface charge density $\sigma_s$ is, by definition, whatever value
makes $\En=0$ hold. It is the system's \emph{response}, not an independent datum.
Imposing the homogeneous Neumann condition \emph{is} the statement ``the equilibrium
surface charge has built up.'' Consequently $\sigma_s$ is never required to
\emph{solve} the field. If one wants it, it is recovered \emph{after} the solve from
Gauss's law at the interface, evaluated just outside the excluded region:
\begin{equation}
  \sigma_s \;=\; \varepsilon_0\,\varepsilon_{\mathrm{LAr}}\;
               \En\big|_{\text{LAr side}} .
  \label{eq:gauss}
\end{equation}
The boundary condition and the surface charge are two descriptions of the same
equilibrium; fixing one fixes the other. We choose to impose $\En=0$ because it is
the quantity the field solver needs and it requires no material constant.

\subsection{Bound (polarization) surface charge is the wrong regime}
The alternative is to carry the \emph{bound} surface charge
$\sigma_b = \mathbf{P}\cdot\nhat$ of a dielectric, which is exactly the displacement
formulation
\begin{equation}
  \nabla\!\cdot\!\big(\varepsilon\,\mathbf{E}\big) = 0 .
  \label{eq:displacement}
\end{equation}
A finite $\varepsilon$ only \emph{refracts} field lines across the interface
(Snell-type bending of $\mathbf{E}$); it never \emph{terminates} them, because a
pure dielectric conducts no charge and never reaches a charged-up steady state. In
the \textsc{pochoir} tests this made the overshoot \emph{worse}: with the
$\varepsilon$ slab the paths reached $z_{\min}=\SI{9.573}{mm}$ with 4 paths pulled
below the pad plane, versus $z_{\min}=\SI{10.000}{mm}$ and 0 paths below with the
no-flux mask. The physical reason is that the free accumulated surface charge, not
the polarization response, governs the boundary of a real insulator held in a DC
field for long times.

\section{What is actually approximated}

Stated honestly, the model makes three approximations, all well justified here:
\begin{enumerate}
  \item \textbf{$\sigma_{\mathrm{FR4}}\to 0$} (perfect insulator): FR4 leakage over
        drift-scale times is negligible given $\rho\sim\SI{e14}{\ohm\,m}$.
  \item \textbf{Steady state}: the one-time transient charging phase is skipped;
        the operating equilibrium is modelled directly.
  \item \textbf{Dielectric refraction dropped}: deliberate --- the free surface
        charge dominates the boundary, so carrying $\varepsilon$ would misdescribe
        the regime (and, as shown, worsens the result).
\end{enumerate}

\section{Comparison of candidate boundary conditions}

\begin{table}[h]
\centering
\begin{tabular}{@{}p{0.30\linewidth} p{0.42\linewidth} c@{}}
\toprule
\textbf{Model} & \textbf{Behaviour at the FR4} & \textbf{Verdict} \\
\midrule
Dirichlet, grounded gap & forces charge collection in the gap
  & non-physical (electron accumulation) \\
Drift-terminate only, field unchanged & $\mathbf{E}$ still points into FR4
  & physically inconsistent \\
Dielectric $\varepsilon$ (displacement, Eq.~\ref{eq:displacement})
  & refracts field lines, never terminates them
  & overshoot persists / worsens \\
\textbf{No-flux Neumann} (Eq.~\ref{eq:neumann})
  & field tangential; charge funnelled to pads
  & \textbf{correct} (charged-up insulator) \\
\bottomrule
\end{tabular}
\caption{Candidate FR4 boundary conditions. Only the homogeneous Neumann condition
reproduces the charged-up-insulator steady state; it needs neither a permittivity
nor an explicit surface-charge density.}
\end{table}

\section{Summary}

The FR4 laminate is modelled as a no-flux (homogeneous Neumann) boundary,
$\partial\phi/\partial n = 0$, equivalent to $\En=0$ at the LAr/FR4 interface. This
is the correct steady-state description of a bulk insulator that cannot conduct the
drift current and therefore charges up until the normal field vanishes. The
accumulated surface charge is not omitted from the physics --- it is the very
response that enforces the boundary condition and is recoverable \emph{a posteriori}
from Gauss's law, Eq.~\eqref{eq:gauss}. Carrying the FR4 as a dielectric
permittivity would model the wrong regime (polarization refraction rather than free
charge accumulation) and empirically worsens the gap overshoot. The no-flux model
introduces no material permittivity and no surface-charge source term, and it
funnels gap electrons onto the pads, eliminating the overshoot.

\end{document}
